% Part II, instruction chapters. The entries are generated by gen/geninsns.py.

\chapter{Integer computation}\label{ch:i-int}
Integer instructions operate on the 64-bit registers. There are no condition flags, no traps on
overflow and no traps on division by zero. All register--register operations share opcode ALU
(\hex{00}) and are selected by \field{funct}; the immediate forms each have their own opcode.
\input{gen-summary-int}
\input{gen-insns-int}

\chapter{Local Store access}\label{ch:i-ls}
Loads and stores move data between the integer registers and the data Local Store, and nothing else:
there is no instruction that loads from or stores to far memory. Addresses are \reg{rs1} plus a
sign-extended 16-bit offset; accesses must be naturally aligned. Every load and store obeys the hazard
rules of Section~\ref{sec:asif}.
\input{gen-summary-ls}
\input{gen-insns-ls}

\chapter{Control transfer}\label{ch:i-ctl}
Branches compare two registers and jump PC-relative. \insn{jal} jumps PC-relative and \insn{jalr}
jumps to a register; both write a return address. There are no delay slots and no condition codes.
\input{gen-summary-ctl}
\input{gen-insns-ctl}

\chapter{Transfers and far atomics}\label{ch:i-xfer}
These instructions are the only way to reach T2, T3 and the devices (Chapter~\ref{ch:memory}), and
\insn{tput} through a core-control capability is the only way to stop another core. Transfers name a
Local Store range, a capability, a far offset, a length and a tag; atomics name a capability and an
offset. \insn{twait}, \insn{tpoll} and \insn{tany} observe tags; \insn{tany} is also the instruction
with which a core waits for events (Section~\ref{sec:events}).
\input{gen-summary-xfer}
\input{gen-insns-xfer}

\chapter{Capability instructions}\label{ch:i-cap}
The capability instructions occupy two opcodes, CAP (\hex{34}) and CAP2 (\hex{35}), and are selected
by \field{funct}\texttt{<3:0>}. CAP holds the operations on segments and inboxes; CAP2 holds revocation
domains (\insn{cdom}, \insn{crevoke}, \insn{cundom}) and endpoints (\insn{csend}, \insn{crecv},
\insn{cmkep}). The rules they implement are described in Chapters~\ref{ch:caps} and \ref{ch:intercore}.
None may overwrite a live linear capability (Section~\ref{sec:nooverwrite}). \insn{csave},
\insn{crestore}, \insn{cdom}, \insn{crevoke}, \insn{cundom} and \insn{cmkep} are privileged and trap with
\cause{priv} in user mode.
\input{gen-summary-cap}
\input{gen-insns-cap}

\chapter{System, cache and overlay instructions}\label{ch:i-sys}
CSR access, halting and trap return (Chapter~\ref{ch:traps}), the software-cache assists
(Section~\ref{sec:swcache}) and the overlay map (Section~\ref{sec:overlays}).
\input{gen-summary-sys}
\input{gen-insns-sys}
